\chapter{Introducción}

\section{Objetivos de la práctica}
El objetivo fundamental de esta práctica es la obtenención de información y 
familiarización con las amenazas a la seguridad digital.

\section{Fundamentos teóricos}
Para la realización de esta práctica, se aplican diversos conceptos teóricos clave 
de la ciberseguridad y la seguridad de la información, entre los que destacan:
\begin{itemize}
    \item \textbf{Conocimiento de los CVE/CWE:} CVE (Common Vulnerabilities and Exposures) es un sistema de identificación de vulnerabilidades conocido a nivel mundial, mientras que CWE (Common Weakness Enumeration) clasifica las debilidades de software. Familiarizarse con estas bases de datos es esencial para entender las amenazas actuales y cómo mitigarlas.
    \item \textbf{Phishing:} Técnica de ingeniería social utilizada para engañar a los usuarios y obtener información confidencial, como contraseñas o datos bancarios, a través de correos electrónicos, mensajes o sitios web falsos.
    \item \textbf{Malware:} Software malicioso diseñado para dañar, explotar o comprometer sistemas informáticos. Puede incluir virus, troyanos, ransomware, spyware, entre otros.
    \item \textbf{Ransomware:} Tipo de malware que cifra los archivos de la víctima y exige un rescate para restaurar el acceso a los datos. Es una amenaza creciente para empresas y particulares.
    \item \textbf{Botnets:} Redes de dispositivos comprometidos (ordenadores zombies) controlados por un atacante para realizar actividades maliciosas, como ataques DDoS, envío de spam o minería de criptomonedas.
    \item \textbf{SQL Injection:} Vulnerabilidad que permite a un atacante ejecutar código SQL malicioso en una base de datos a través de una aplicación web, lo que puede resultar en la exposición o manipulación de datos sensibles.
    \item \textbf{Cross-Site Scripting (XSS):} Vulnerabilidad que permite a un atacante inyectar scripts maliciosos en páginas web vistas por otros usuarios, lo que puede llevar al robo de cookies, secuest
\end{itemize}

\section{Enunciado de la práctica}
Esta práctica consiste en la realización de un informe basado en una selección de noticias de prensa digital, posteriores al 1 de enero de 2025, relacionadas con una serie de amenazas informáticas.

Las amenazas objeto de esta actividad serán:
\begin{enumerate}
    \item Phishing, estafas a través de la red
    \item Malware difundido en redes sociales
    \item Ataques de ransomware contra corporaciones y gobiernos
    \item Botnets, ordenadores zombies
    \item Malware en dispositivos móviles
    \item Inyección de código SQL
    \item Cross-Site Scripting (XSS)
\end{enumerate}

